\bookchapter{Spin across substrate clocks, matter and exact computation}
\label{ch:spin-connections}

\subsection{One research chain, with explicit receiving maps}

The completed spin program connects the three volumes through concrete mathematical objects. Volume I supplies source, clock, propagation and physical-record questions. The spin section of Volume II supplies complete interacting preparations, hidden-state response, retained histories and calculated receiving bounds. The following matter chapters receive those methods through actual isotope, atomic and molecular source maps. Volume III supplies exact factorization, native operator composition and persistent computational objects. This chapter organizes those connections as an editorial synthesis of the recorded interfaces.\samref{SPX_VOLI}\samref{SPX_VOLIII}

The useful chain is
\begin{gather}
 \begin{gathered}
 \text{source action and preparation}\\
 \longrightarrow\text{ordered response},
 \end{gathered}\\
 \begin{gathered}
 \text{ordered response}\longrightarrow\text{paid receiver}\\
 \longrightarrow\text{retained computational object}.
 \end{gathered}
\end{gather}
The object passed along this chain includes the source identity, normalization, state or Gibbs law, actual generator, clock and receiving observable. A completed thermal factor can be reused by another recipient, but its numerical coefficient still belongs to the source that generated it. These mappings make the connection constructive: the next calculation receives identified operands rather than a shared label.

\begin{table*}[tbp]\centering\small
\begin{tabularx}{\textwidth}{@{}l*{3}{>{\raggedright\arraybackslash}X}@{}}\toprule
Interface & Volume I & Volume II & Volume III\\\midrule
Source & Energy, field and reference & Complete spin or matter preparation & Exact source object and factors\\
Clock & Proper age and source phase & Original attempt law and memory & Generator, elapsed operand and ordered calls\\
History & Propagation and retained record & Post-contact state and hidden response & Event order, dependency graph and checkpoint\\
Readout & Detector, calibration and work & Thermal, energy and response bounds & Certified matrices, contractions and recovery\\\bottomrule
\end{tabularx}
\end{table*}

The entries identify a division of work within one research program. The source-specific results in the preceding chapters supply the equations and evidence for the links.

\subsection{The clock connection is a statement about the full generator}

Volume I retains the static local lift and endpoint lapse
\begin{equation}
 A_{\rm lift}(r)=\frac{r_s}{r},\qquad
 \frac{d\tau}{dt}=\sqrt{1-A_{\rm lift}},
 \qquad r_s=\frac{2GM}{c^2}.
\end{equation}
Its action-phase convention is $\Delta\phi=S/\hbar$. The formation floor $A_0=1/(12\pi)$ has its own shared-reference normalization. The global accumulated road $A_{\rm los}=\pi^{-1}[1-(1+z)^{-3}]$ has a different source and observational role from a local lift. These physical observables provide specified recipients for a microscopic construction.\samref{SPX_VOLI}

The exact spin clock criterion identifies the pure common-rate sector. If every boundary and hidden attempt receives the same positive multiplier, with dimensionless energy gaps fixed, then
\begin{equation}
 L_A=\ell(A)L_0.
\end{equation}
The Gibbs law, conditional hidden weights and rate-overlap structure remain unchanged. For a time-dependent common multiplier the same full-state protocol uses operational time
\begin{equation}
 \tau(t)=\int_0^t\ell(A(u))\,du.
\end{equation}
Matched $\tau$ compares the same ordered generator evolution; fixed laboratory waits can distinguish the resulting histories. The retained assay uses the declared example $\ell(A)=\sqrt{1-A}$ with lifts zero and $1/4$. It states the source-lift assignment and scales all attempts together.\samref{SPX_CHANNEL}

Changing source composition or dimensionless gaps changes another part of the dynamics. Chapter~\ref{ch:spin-composition} calculates the exact source-dependent overlap between microscopic and conditional thermal rates. Its edge dependence and memory are retained in addition to any common kinetic multiplier. The later source-clock construction uses a separate compiled seven-spin, 128-state gate, with one output bit, six hidden spins, inverse temperature one and unit attempts. Write $B(s)$ for the sum of its hidden signs divided by six. For its literal strain $E_A=E_0-AB$, bank activity ratios have opposite signs relative to baseline: two increase and two decrease for every $0<A\leq1$. A single positive common prefactor cannot turn these four distinct shape responses into one prescribed activity ratio.\samref{SPX_CLOCK}\samref{SPX_CLOCK_PROOF}

This supplies a practical physical connection. A proposed source-to-lapse map must carry both the common kinetic or phase factor and the actual source-shape response. In the retained canonical-field example, put $z_s=1+B(s)/2$, inertia $\mu_A(s)=1+Az_s$ and restoring coefficient $\chi=1$. Its frequency factor $(1+Az_s)^{-1/2}$ and the static lapse $\sqrt{1-Az_s}$ share their weak half slope but differ at second order. The exact lapse requires the constitutive ratio $\chi/\mu=1-Az_s$. The source calculation identifies the response that this physical assignment must receive.\samref{SPX_CLOCK_PROOF}

The spin histories in the latest response chapter keep their original $N/\xi$ uniform site-attempt clock. Their certified intervals are available in that clock. A physical seconds map is an additional source/action assignment, with its retained source normalization, rather than an adjustment of those completed intervals.\samref{SPX_DIRICHLET}

\subsection{The field and record connection retains chronology}

The spin memory mechanism explains why a present coarse reading can leave out information that affects the next interaction. Chapter~\ref{ch:spin-microscopic} retains the hidden-state forcing, memory kernel, preparation and post-reading state. The physical response chapter then uses source functions and their full-Gibbs norms to obtain paid bounds on a collective receiver. These results distinguish local conditional relaxation from a later full-source response.

Volume I's current shared-field construction asks a related question with an explicit continuum quantum source and actual event propagators. Two prescribed source worldtubes act in one common Fock field. Its comoving generator is
\begin{equation}
 K_j=H_f-v_j\cdot P_f+n_j(H_0-H_f),
 \qquad n_j=\sqrt{1-|v_j|^2}.
\end{equation}
Translation mixes field modes, and the source-driven field retains its interacting currents and backreaction. For physically specified histories, the record contains the coherence matrix
\begin{equation}
 G_{wv}=\operatorname{Tr}(\rho U_v^*U_w),
\end{equation}
which is positive with unit diagonal. Equal endpoints or action totals do not determine this ordered operator matrix. Actual accessible conditional states determine the decoding instrument. The completed chapter carries route, encoder, decoder, detector and retained-memory costs through the same pipeline.\samref{SPX_SHARED}

The connection to spin is a precise research method: keep the full state selected by an earlier interaction when calculating a later response. The classical Gibbs/heat-bath memory source and the continuum unitary source keep their own generators. The shared-field report explicitly records that no spin coefficients enter its receiving model. Its mathematical use of full mixed terms, physical background and retained correlations is the transferable lesson. The new spin section now gives that lesson its original source derivation and completed response evidence.

The response estimates likewise carry different kinds of payment. A small state-distance error controls a bounded final observable through its receiving channel. An energy or action norm requires its own source moments and field/background bounds. The latest spin Dirichlet result supplies one such source-specific paid moment. The shared-field construction separately carries actual encoder and switching energy. Both preserve the relationship between a particular error estimate and the observable it bounds.\samref{SPX_DIRICHLET}\samref{SPX_SHARED}

\subsection{The matter connection uses the actual nuclear action}

The link to the following matter chapters is already executable. For the retained Li6 source, nine contact phases $z_e\in\{1,i,-1,-i\}$ have the exact binary chart
\begin{equation}
 z_e=\frac{s_e+t_e}{2}+i\frac{s_e-t_e}{2},
 \qquad s_e,t_e\in\{-1,+1\}.
\end{equation}
The full source quadratic form becomes
\begin{equation}
 S_m=\sigma^TG_m\sigma
 =\operatorname{tr}G_m-\sum_{i<j}J_{ij}\sigma_i\sigma_j,
 \qquad J_{ij}=-2(G_m)_{ij}.
\end{equation}
All self, motif and hidden-loop terms remain in the complete action. The eighteen binary coordinates encode the nine original phases while the material inventory stays $3p3n$. The original source result explicitly retains that distinction. It calculates the four physical join/separation masks and their connected response on the same phase addresses.\samref{NUC_ISOTOPE}

The hidden-spin quartet module supplies a constructive ground interaction for the occupied contact interface. Its absolute reference and degeneracies then matter to a finite-temperature recipient. The completed comparison finds that the original isotope and ground-matched auxiliary module have different thermal responses at the same source beta. The matter construction therefore receives the actual original conditional spectra. It keeps 256 boundary rows with 1024 hidden states each, every source constant and the original Gibbs preparation.\samref{NUC_THERMAL}

For that preparation the same-address work and free-energy change satisfy
\begin{equation}
 D(p_{00}\Vert p_m)=\beta\big[W_m-(F_m-F_{00})\big].
\end{equation}
This is the source-specific place where hidden relaxation, thermal interaction and payment enter a nuclear contact. It is also why the quartet's exact ground realization and its different thermal factor both belong in the book. The completed thermal stage supplies the correct factors for larger nuclear assemblies and electronic recipients.

The later isotope histories then retain the changed wave, actual contact coordinate, selected incoming preparation and full complete-energy difference. The occupied atom receives that history through its moving Pauli propagators and conditional channel. The molecule receives the same source with its full moving-fragment reference. These connections are developed in the following chapters using their actual coefficients. The spin source supplies a method and, where the explicit chart permits it, exact operands; each physical recipient retains its source-to-field and energy map.

\subsection{A common energy family is an achieved intermediate connection}

The completed interaction-metric work maps retained source interactions into two preparation-independent energy families on the original 54 spin preparations. Its actual current incidence has rank fourteen and leaves four tensor/scalar parity modes to complete. The full compatible parity-energy kernel is
\begin{equation}
 \alpha P_{TT}+\beta_LP_L+\beta_RP_R.
\end{equation}
The coefficients specify one tensor-pair energy and separate scalar-parity energies. They are common across preparations. A second, literal two-object incidence retains all sixteen port modes and its signed-turn work.\samref{SPX_METRIC}

The calculation carries the coherent mean, fluctuations, six-function forcing, constant and higher complement, and their signed interference. The recorded matrix identities match the six-coordinate energy calculation with the complete eighteen-coordinate source form on all 54 preparations. The older covariance-dual metrics remain available with their preparation dependence. The Hamiltonian forms have a separate physical role.

This result gives the cross-volume source question a finite list of inputs: actual role/current incidence, the common interaction and parity coefficients, field-force normalization, receiving scale, and the action/radius/seconds map. The same campaign retains a complete receiving norm conversion, including acceptance, failure and arbitrary ancillas. Its spin packet and completed electron controls have different preparations; the common number of controls supplies no isotope row pairing. The achieved connection is the explicit energy-and-receiver family. Its remaining physical selections can be carried into Volume I's source field and Volume II's matter energies without changing the recorded mathematical results.\samref{SPX_METRIC}

\subsection{Volume III receives exact states and reusable operations}

The computational companion receives several different finished objects. A joint density supplies every integer multiplicity $g(E,M,b)$. A compact thermal source supplies the exact factor definition for a specified directed integral. A dynamical source supplies its actual generator, conditional weights, initial state and readout. A certified receiver supplies matrix intervals and their physical norm payments. A retained native object supplies the bound mathematical value together with the source and dependency closure needed for later composition.\samref{SPX_VOLIII}

Exact factorization preserves the complete mathematical state space. The packet construction evaluates local polynomials and signed bridges while retaining the dense interaction through magnetization. Permutation grouping retains binomial multiplicities for every hidden class. Krylov or receiving-coordinate calculations retain or bound their complementary source action. These are distinct ways of reducing work with the required source information still present.

The shared joint-spin installation already adopts those sources in Starbreaker and ATOM3D. Later recorded spin methods compose source projection, hidden residual maps, residual payment, generalized directions and receiving lower bounds through native retained references. The installed source-clock route also receives certified signed-Write integrals as elapsed operands. Its original trajectory results and qualification records retain their own execution versions. The five recorded spin compositions are a concrete Volume III interface for the mathematical chain developed in this section.\samref{SPX_INSTALL}

A representative receiving dependency is
\begin{equation}
 \begin{gathered}
 (G_5,D_5,\text{source moments})\\
 \longrightarrow\text{projected residual}\\
 \longrightarrow\text{paid moment}\\
 \longrightarrow\text{Schur response bound}.
 \end{gathered}
\end{equation}
The source, matrix and payment dependencies survive export and fresh recovery. A later consumer can receive the completed operand without repeating its source integral. The current Dirichlet result uses exactly this separation: retained normalizers and source moments combine with a newly evaluated edge observable, which then feeds the six-coordinate receiver. The computation is part of the scientific construction because it preserves which actual source supplied each operand.\samref{SPX_DIRICHLET}

The more recent layered dense-source solver adds another completed computational route. It verifies an integer decomposition
\begin{equation}
 J=\sum_k a_k(v_kv_k^T-I)+S,
\end{equation}
then carries layer magnetizations through sparse residual elimination and an exact min-plus join. The floating calculation only proposes candidate layer vectors; the decomposition and final accounting are checked exactly. Its output is an exact ground energy, degeneracy and witness with exhaustive $2^N$ accounting, rather than a full joint density by the same name. In the retained qualification, 52 of 60 planted small sources admit and agree with complete enumeration. Both later raw dense-source solve tasks completed; their separate selection tasks failed on an absent mathematical object store. The report preserves both outcomes.\samref{SPX_LAYERED}

This computational addition belongs naturally in Volume III's route selection and exactness account. In Volume II it identifies the mathematical object available to the matter researcher. A ground-state certificate, complete density, thermal trace and time-dependent response each remain available for their own receiving questions.

\subsection{The completed bridge and the next source operands}

The dedicated spin section now supplies a continuous scientific account: exact state organization, finite thermal structure, hidden memory, composition control, actual kinetic response, physical source moments and paid receiving bounds. Its connections to matter are explicit phase/action and conditional-preparation maps. Its connections to Volume I identify the full generator and physical source/clock/field recipients. Its connections to Volume III identify the exact data, operators and retained dependencies that make the research reusable.

The latest source review retains a completed one-residual Dirichlet upper and a separate six-coordinate lower receiver. The full mixed two-residual Dirichlet block and next squared-action norm remain the named source targets. C2 optimization remains stopped/incomplete and the RAM comparison remains partial in their own records. These statuses identify the next mathematical operands without turning them into results already present. The following matter chapters can use the completed source constructions now, with their physical coefficients, references and units carried explicitly.\samref{SPX_DIRICHLET}
